\section{Combinatorial mutations of perfect matching polygons}\label{sec_mutation_polygon1}
In this section, we discuss a relationship between deformations of consistent dimer models and combinatorial mutations of polygons.
We first define combinatorial mutations for lattice polytopes of any dimension, and then we mainly discuss the case of polygons.

\subsection{Preliminaries on combinatorial mutations of polytopes}\label{sec_pre_mutation_polygon}

Following \cite{ACGK}, we introduce combinatorial mutations of polytopes.
Let $N\cong\ZZ^d$ be a lattice of rank $d$, and $P\subset N_\RR\coloneqq N\otimes_\ZZ\RR$ be a convex lattice polytope.
We assume that $P$ contains the origin ${\bf 0}$.
We denote by $\calV(P)$ the set of vertices of $P$.
We say that two polytopes $P,Q\subset N_\RR$ are \emph{isomorphic} if they are transformed into each other by $\GL(d,\ZZ)$-transformations, in which case we denote $P\cong Q$.

We first prepare some notions used in the definition of combinatorial mutations of polytopes.

\begin{Definition}[mutation data]\label{def_mutation_data}
Let $P$ be a convex lattice polytope as above.
In order to define a combinatorial mutation of $P$, we fix the data consisting of a vector $w$ and the associated integers $h_{\max}(P,w)$,
$h_{\min}(P,w)$ and $\width(P,w)$ defined as follows.

Let $w\in M\coloneqq\Hom_\ZZ(N,\ZZ)\cong\ZZ^d$ be a primitive lattice vector.
The element $w\in M$ determines the linear map $\langle w,-\rangle\colon N_\RR\rightarrow \RR$, where $\langle-,- \rangle$ is the natural inner product.
We set
\begin{displaymath}
h_{\max}(P,w)\coloneqq{\max}\{\langle w,u\rangle \,|\, u\in P\} \qquad\text{and}\qquad
h_{\min}(P,w)\coloneqq{\min}\{\langle w,u\rangle \,|\, u\in P\},
\end{displaymath}
which are integers since $P$ is a lattice polytope.
When the situation is clear, we simply denote these by $h_{\max}$ and $h_{\min}$, respectively.
We define the \emph{width} of $P$ with respect to $w$ as
\begin{displaymath}
\width(P,w)\coloneqq h_{\max}(P,w)-h_{\min}(P,w).
\end{displaymath}
\end{Definition}

We note that if the origin ${\bf 0}$ is contained in the strict interior $P^\circ$ of $P$, then $h_{\min}<0$ and $h_{\max}>0$, in which case we have $\width(P,w)\ge2$.
We say that a lattice point $u\in N$ (resp.\ a~subset $F\in N_\RR$) is at \emph{height} $m$ with respect to $w$ if $\langle w,u\rangle=m$
(resp.\ $\langle w,u\rangle=m$ for any $u\in F$).

For each height $h\in\ZZ$, we let
\begin{displaymath}
w_h(P)\coloneqq\operatorname{conv}\{u\in P\cap N \,|\, \langle w,u\rangle=h\},
\end{displaymath}
which is the (possibly empty) convex hull of all lattice points in $P$ at height $h$.
By definition, $w_{h_{\min}}(P)$ and $w_{h_{\max}}(P)$ are faces of $P$.
Using these notions, a combinatorial mutation of $P$ is defined as follows.

\begin{Definition}\label{def_mutation_polygon}
Let the notation be the same as above.
We assume that there exists a lattice polytope $F\subset N_\RR$ such that $\langle w,u\rangle=0$ for any $u\in F$, and
for each negative height $h_{\min}\le h<0$ there exists a possibly empty lattice polytope $G_h\subset N_\RR$ satisfying
\begin{equation}\label{condition_Gh}
\{u\in\calV(P)\,|\, \langle w,u \rangle=h\}\subseteq G_h+(-h)F\subseteq w_h(P),
\end{equation}
where $+$ means the Minkowski sum, and we especially define $Q+\varnothing=\varnothing$ for any polytope $Q$.
We call $F$ a \emph{factor} of $P$ with respect to~$w$.
Then, we define the \emph{combinatorial mutation} of $P$ given by the vector $w$, factor $F$ and polytopes $\{G_h\}$ as
\begin{displaymath}
\mutation_w(P,F)\coloneqq\operatorname{conv}\bigg(\bigcup_{h=h_{\min}}^{-1}G_h\cup\bigcup_{h=0}^{h_{\max}}(w_h(P)+hF)\bigg).
\end{displaymath}
\end{Definition}

We note that the combinatorial mutation is independent of the choice of the polytopes $\{G_h\}$ (see \cite[Proposition~1]{ACGK}).
Also, a translation of the factor $F$ does not affect the combinatorial mutation; that is,
for any $u\in N$ with $\langle w,u\rangle=0$ we have $\mutation_w(P,F)\cong\mutation_w(P,u+F)$; see~\cite{ACGK} for more details.

\begin{Remark}[the combinatorial mutation for the case of $d=2$]\label{rem_def_mutation}
When $d=2$, we choose an edge $E$ of a lattice polygon~$P$, and take $w\in M\cong\ZZ^2$ as a primitive inner normal vector for~$E$.
By the choice of $w$, we see that $w_{h_{\min}}(P)=E$ and $w_{h_{\max}}(P)$ is either a vertex or an edge of~$P$.
Then, we take a primitive lattice element $u_E\in N$ satisfying $\langle w,u_E \rangle=0$, and define the line segment $F\coloneqq\operatorname{conv}\{{\bf 0},u_E\}$,
which is parallel to $E$ at height $0$ and has unit lattice length. Since $u_E$ is uniquely determined up to sign, so is $F$.
In this case, $P$ admits a combinatorial mutation with respect to $w$ (equivalently, there exit polytopes $\{G_h\}$ satisfying (\ref{condition_Gh})) if and only if $|E\cap N|-1\ge -h_{\min}$, see \cite[Lemma~1]{KNP}.
We note that the combinatorial mutation does not depend on the choice of $u_E$ (and hence $F$). That is,
$\mutation_w(P,F)\cong\mutation_w(P,-F)$, which means they are $\GL(2,\ZZ)$-equivalent.
\end{Remark}
